% ============================================================
% Section 181: 近简并2s/2p能级在电场中的行为
% ============================================================
\section{量子力学：近简并\texorpdfstring{$2s$}{2s}/\texorpdfstring{$2p$}{2p}能级在电场中的行为}
\label{sec:near-degenerate-2s2p}

\begin{modelbox}[题目：近简并\texorpdfstring{$2s$}{2s}/\texorpdfstring{$2p$}{2p}能级在电场中的行为]
假定上题中 $|210\rangle$ 和 $|200\rangle$ 不简并，能量分别为 $E^0+\Delta$ 和 $E^0-\Delta$。
\begin{enumerate}
    \item 证明 $H_0+H'$ 关于这两态的矩阵具有本征值 $E^0\pm\sqrt{\Delta^2+A^2}$，$A=3ea_0\mathcal{E}$；
    \item 将极限 (i) $A\gg\Delta$ 和 (ii) $A\ll\Delta$ 与微扰论结果比较；
    \item 画出 $E^0\pm\sqrt{A^2+\Delta^2}$ 随 $\mathcal{E}$ 的函数关系；
    \item 若波数差为 $36\,\text{m}^{-1}$，计算使 $A=\Delta$ 的电场强度。
\end{enumerate}
\end{modelbox}

\begin{tipbox}[考点快速分析]
\begin{itemize}
    \item \textbf{二能级系统}：精确对角化 $2\times2$ 矩阵，避免交叉
    \item \textbf{极限对比}：$A\gg\Delta$ 对应简并微扰（一级线性Stark）；$A\ll\Delta$ 对应非简并微扰（二级平方Stark）
    \item \textbf{能级反交叉}：$E(\mathcal{E})$ 曲线永不相交，反映无交叉定理
    \item \textbf{数值估算}：波数 $\to$ 能量 $\to$ 电场强度
\end{itemize}
\end{tipbox}

\begin{derivationbox}[标准解题过程]
\textbf{1. 矩阵对角化}

在 $\{|210\rangle,|200\rangle\}$ 基中：
\[
H=\begin{pmatrix}E^0+\Delta&-A\\-A&E^0-\Delta\end{pmatrix}.
\]

久期方程：$(E^0+\Delta-\lambda)(E^0-\Delta-\lambda)-A^2=0$，即 $(E^0-\lambda)^2-\Delta^2-A^2=0$。

\begin{equation}
\boxed{\lambda_\pm=E^0\pm\sqrt{\Delta^2+A^2}.}
\end{equation}

\textbf{2. 极限与微扰论比较}

\textit{(i) $A\gg\Delta$（强场/近简并）}：
\[
\lambda_\pm\approx E^0\pm A\pm\frac{\Delta^2}{2A}.
\]
领头项 $E^0\pm A$ 正是简并微扰论的一级修正（题 7.47）。

\textit{(ii) $A\ll\Delta$（弱场/非简并）}：
\[
\lambda_\pm\approx E^0\pm\Delta\pm\frac{A^2}{2\Delta}.
\]
$\pm A^2/(2\Delta)$ 对应非简并微扰论的二级修正（分母为能级差 $2\Delta$）。

\textbf{3. 能级图}

以 $\mathcal{E}$（即 $A$）为横坐标：
\begin{itemize}
    \item $\mathcal{E}=0$：两曲线交于 $E^0\pm\Delta$，对应纯 $|210\rangle$（上）和 $|200\rangle$（下）。
    \item $A\gg\Delta$：渐近于 $E^0\pm A$（直线），本征态趋于对称/反对称组合。
\end{itemize}
曲线呈双曲形状，永不相交——\textbf{避免交叉}（avoided crossing）。

\textbf{4. 数值计算}

波数差 $\Delta\tilde{\nu}=36\,\text{m}^{-1}$，能量差 $2\Delta=hc\Delta\tilde{\nu}$。
\[
\Delta=\frac{hc\Delta\tilde{\nu}}{2}\approx\frac{1.986\times10^{-25}\times36}{2}\approx3.57\times10^{-24}\,\text{J}.
\]

令 $A=3ea_0\mathcal{E}=\Delta$：
\begin{equation}
\boxed{\mathcal{E}=\frac{\Delta}{3ea_0}\approx\frac{3.57\times10^{-24}}{3\times1.602\times10^{-19}\times0.529\times10^{-10}}\approx1.4\times10^5\,\text{V/m}.}
\end{equation}
\end{derivationbox}

\begin{formulabox}[答案汇总]
\begin{center}
\begin{tabular}{ll}
\toprule
\textbf{结论} & \textbf{结果} \\
\midrule
本征值 & $\lambda_\pm=E^0\pm\sqrt{\Delta^2+A^2}$ \\
强场极限 & $\lambda_\pm\approx E^0\pm A$（简并微扰） \\
弱场极限 & $\lambda_\pm\approx E^0\pm\Delta\pm A^2/(2\Delta)$（非简并二级） \\
$A=\Delta$ 电场 & $\approx1.4\times10^5\,\text{V/m}$ \\
\bottomrule
\end{tabular}
\end{center}
\end{formulabox}

\begin{insightbox}[物理图像]
\begin{itemize}
    \item 避免交叉是量子力学中最普遍的定理之一：具有相同对称性的两个能级在外参数变化时不会相交，最小能隙正比于耦合强度（此处为 $A$）。
    \item 实际氢原子中，$2s_{1/2}$ 与 $2p_{1/2}$ 因 Lamb 位移存在约 $1058\,\text{MHz}$ 的能差（对应 $\Delta\tilde{\nu}\approx1.2\,\text{m}^{-1}$），远小于本题的 $36\,\text{m}^{-1}$。这使氢原子在弱电场下即可表现出线性 Stark 效应。
    \item 该模型也是理解量子比特和分子光谱中 ``avoided crossing'' 现象的基石。
\end{itemize}
\end{insightbox}
